\documentclass[10pt,a4paper]{amsart}
\usepackage[margin=19mm]{geometry}
\usepackage{amsmath,amssymb,mathtools}
\usepackage[scr=rsfs,cal=euler]{mathalfa}
\usepackage{xfrac}
\usepackage[unicode,hidelinks,bookmarks=false]{hyperref}
\newcommand{\ie}{i.e.\ }
\newcommand{\cf}{cf.\ }
\newcommand{\resp}{respectively\ }
\newcommand{\Subsection}{\subsection}
\providecommand{\nobreakdash}{\nobreak\hbox{-}}
\setlength{\emergencystretch}{2em}
\multlinegap0pt
\usepackage{amssymb}
\usepackage{natbib}
\usepackage{booktabs}
\usepackage[shortlabels]{enumitem}
\usepackage{xcolor}
\usepackage{algorithm}\usepackage{algpseudocode}
\usepackage{caption}
\newtheorem{proposition}{Proposition}
\title{A generalized Hurwitz stability criterion via rectangular block Hankel matrices for nonmonic matrix polynomials}
\author{Zixiang Ni, Yongjian Hu, Xuzhou Zhan}
\date{Source: arXiv:2508.14376v7 (CC BY 4.0)}
\begin{document}
\maketitle

\noindent\emph{Source and licence.} A generalized Hurwitz stability criterion via rectangular block Hankel matrices for nonmonic matrix polynomials. Version arXiv:2508.14376v7 (CC BY 4.0), distributed by arXiv under the Creative Commons Attribution 4.0 International licence, http://creativecommons.org/licenses/by/4.0/, as recorded in the arXiv licence metadata for this exact version and shown on the abstract page https://arxiv.org/abs/2508.14376v7. Full licence notice: \texttt{licenses/arxiv-2508.14376-CC-BY-4.0.txt}.\par\smallskip

\noindent\textit{Editorial scope.} Counted unit: the article's proposition LemHs (source0.tex lines 492--499). The article's complete original proof of it is delivered with it (source0.tex lines 501--537, one \texttt{proof} environment of 37 lines). No mathematics is written, repaired or re-derived: the statement and the proof are the article's own lines, in the article's own notation, and no retained line is substituted or re-flowed.  Retained uncounted as the source-local context the proof needs: lines 338--340; lines 346--348.  Nothing was omitted from any retained span, so the package carries no omission record.  No retained line contains a citation, so the wrapper carries no bibliography and no bibliography entry.  No retained line contains a figure environment, an image-inclusion command or drawing code, so no image is delivered and the package declares no asset dependency.  Editorial formatting only, in addition: an AMS \texttt{amsart} preamble that restates the article's own declarations and macros used by the retained lines, namely the environments \texttt{proposition} on separate counters, together with the standard packages those lines need.  The retained environments print in this document as Proposition 1 in order of appearance, whereas the article prints the same items with the numbering of its own full document.  The complete native source of the article as distributed in the arXiv e-print for this exact version is bundled unchanged as \texttt{sources/183546/source0.tex}.
\begin{equation}\label{LaurentSeriesRsp}
       R(\lambda)={\bf s}_{-1}+\sum_{j=0}^{\infty} \frac{{\bf s}_j}{\lambda^{j+1}},\quad \lambda\rightarrow \infty.
     \end{equation}

\begin{equation}\label{Matrixfrac}
R(\lambda)=(D_l(\lambda))^{-1}N_l(\lambda)=N_r(\lambda)(D_r(\lambda))^{-1},
\end{equation}

\begin{proposition}\label{LemHs}
Let $R(\lambda)\in \mathbb C(\lambda)^{p\times p}$ be a strictly proper rational matrix that admits the Laurent series \eqref{LaurentSeriesRsp} and matrix fraction decomposition \eqref{Matrixfrac}, where $\{D_l(\lambda), N_l(\lambda),D_r(\lambda), N_r(\lambda)\}\subseteq \mathbb C[\lambda]^{p\times p}$, $D_l(\lambda)$ is row reduced, and $D_r(\lambda)$ is column reduced.
Then \begin{equation}\label{HR}
\mathbf{H}(D_l,N_l;N_r,D_r)=({\bf  s}_{i,j})_{i,j=0}^{\deg D_r-1,\deg D_l-1},
\end{equation}
where ${\bf  s}_{i,j}:=\left({\bf s}_{i+j}\right)_{\mathcal I^c_i(D_r)}^{\mathcal I^r_j(D_l)}\in |\mathcal I^c_i(D_r)|\times |\mathcal I^r_j(D_l)|$.

\end{proposition} 

\begin{proof} Let $k\in \{1,\ldots,p\}$. If ${\rm rdeg}_kD_l=0$,  $[k]^r_{D_l}=\emptyset$. When ${\rm rdeg}_kD_l>0$,  $$[k]^r_{D_l}=\{k,p+k,\ldots,({\rm rdeg}_kD_l-1)p+k\}.$$ In this case,
the columns of $\mathscr{H}_{\infty}( R)^{[k]^r_{D_l}}$ consist of all $k$-th columns of the submatrices 
\begin{equation*}
%\label{column}
\begin{bmatrix}
{\bf s}_{0}\\
{\bf s}_{1}\\
\vdots
\end{bmatrix},\begin{bmatrix}
{\bf s}_{1}\\
{\bf s}_{2}\\
\vdots
\end{bmatrix},\ldots,\begin{bmatrix}
{\bf s}_{{\rm rdeg}_kD_l-1}\\
{\bf s}_{{\rm rdeg}_kD_l}\\
\vdots
\end{bmatrix},
\end{equation*}
Note that $\deg D_l=\max \{{\rm rdeg}_kD_l:k=1,\ldots,p\}$.
To sum up, the columns of $\mathscr{H}_{\infty}(R)^{\mathcal I_r(D_l)}$ contain all $k$-th columns of 
$
\begin{bmatrix}
{\bf s}_{j}\\
{\bf s}_{j+1}\\
\vdots
\end{bmatrix},
$
where $k\in \mathcal I^r_j(D_l)$ and $j=0,\ldots,\deg D_l-1$. Analogously, the rows of $\mathscr{H}_{\infty}(R)_{\mathcal I_c(D_r)}$ contain all $k$-th rows of 
$
\begin{bmatrix}
{\bf s}_{i} &
{\bf s}_{i+1} &\cdots
\end{bmatrix},
$
where $k\in \mathcal I^c_i(D_r)$ and $i=0,\ldots,\deg D_r-1$. Thus, the equation \eqref{HR} holds.

\end{proof}

\end{document}
